\documentclass[14pt,a4paper]{extarticle}

\usepackage{fontspec}
\usepackage{polyglossia}
\setmainlanguage{russian}
\setotherlanguage{english}
\setmainfont{Times New Roman}
\setmonofont{Courier New}
\newfontfamily\cyrillicfont{Times New Roman}
\newfontfamily\cyrillicfonttt{Courier New}

\usepackage[
  left=30mm,
  right=15mm,
  top=20mm,
  bottom=20mm
]{geometry}
\usepackage{setspace}
\usepackage{indentfirst}
\usepackage{titlesec}
\usepackage{graphicx}
\usepackage{float}
\usepackage{caption}
\usepackage{tocloft}
\usepackage{hyperref}

\onehalfspacing
\setlength{\parindent}{1.25cm}
\setlength{\parskip}{0pt}

\titleformat{\section}
  {\normalfont\bfseries\centering}
  {\thesection}
  {1em}
  {\MakeUppercase}
\titleformat{\subsection}
  {\normalfont\bfseries}
  {\thesubsection}
  {1em}
  {}

\captionsetup{labelsep=endash,justification=centering}
\renewcommand{\cftsecleader}{\cftdotfill{\cftdotsep}}
\hypersetup{
  colorlinks=true,
  linkcolor=black,
  citecolor=black,
  urlcolor=black
}

\begin{document}

\begin{titlepage}
  \begin{center}
    Университет ИТМО\\
    Факультет \textless{}название факультета\textgreater{}

    \vfill

    {\bfseries ВЫПУСКНАЯ КВАЛИФИКАЦИОННАЯ РАБОТА}\\[1cm]
    на тему:\\
    {\bfseries \textless{}тема ВКР\textgreater{}}

    \vfill

    \begin{flushright}
      Выполнил: \textless{}ФИО студента\textgreater{}\\
      Группа: \textless{}номер группы\textgreater{}\\
      Руководитель: \textless{}ФИО руководителя\textgreater{}
    \end{flushright}

    \vfill

    Санкт-Петербург\\
    2026
  \end{center}
\end{titlepage}

\tableofcontents
\newpage

\section*{Введение}
\addcontentsline{toc}{section}{Введение}

Современные инструменты проектирования интерьеров позволяют создавать
фотореалистичные сцены, однако значительная часть работы по-прежнему выполняется
вручную. Специалисту необходимо подобрать предметы мебели и отделочные
материалы, расставить их в помещении, настроить освещение, подготовить визуальные
материалы и сформировать спецификацию с товарами и стоимостью. Такой процесс
требует времени, зависит от опыта исполнителя и плохо масштабируется при
необходимости быстро подготовить несколько вариантов оформления одного и того же
помещения.

Актуальность работы связана с необходимостью автоматизации подготовки
дизайн-предложений для помещений на основе исходной геометрии комнаты и каталога
доступных товаров. В практическом сценарии важно не только сгенерировать
визуально правдоподобную расстановку, но и связать элементы сцены с реальными
товарами, их изображениями, ценами и материалами, чтобы результат можно было
использовать в коммерческом процессе.

Целью работы является разработка программного комплекса для автоматизированной
генерации вариантов расстановки предметов интерьера, подбора отделочных
материалов и формирования отчета с визуализациями, товарами и стоимостью.

Для достижения цели были поставлены следующие задачи:
\begin{enumerate}
  \item проанализировать существующие подходы к генерации 3D-сцен и расстановке
  объектов в помещениях;
  \item выбрать базовый инструмент генерации сцен и определить ограничения,
  возникающие при его применении к задаче интерьерного дизайна;
  \item разработать пайплайн улучшения расстановки объектов, настройки света и
  замены синтетических объектов на товары из каталога;
  \item подготовить механизм описания предметов и отделочных материалов с
  использованием VLM- и LLM-моделей;
  \item собрать каталог товаров и материалов, пригодный для автоматической
  подстановки в сцену и расчета стоимости;
  \item провести эксперименты и оценить качество получаемых результатов;
  \item реализовать выгрузку результата в виде изображений, файла сцены,
  спецификации товаров и итогового отчета.
\end{enumerate}

Объектом исследования является процесс подготовки интерьерных 3D-сцен и
дизайн-предложений. Предметом исследования являются методы автоматической
расстановки объектов, подбора материалов и сопоставления синтетических объектов
с товарами из внешнего каталога.

Практическая значимость работы заключается в возможности сократить ручной труд
при подготовке дизайн-предложений и получить воспроизводимый пайплайн, который
может быть интегрирован в инфраструктуру компании.

\section{Аналитический обзор и выбор базового подхода}

\subsection{Автоматизация подготовки интерьерных решений}

В данном разделе нужно описать, почему автоматизация расстановки и подбора
товаров является важной практической задачей. Здесь следует раскрыть следующие
аспекты:
\begin{enumerate}
  \item ручная подготовка интерьера требует большого количества итераций:
  планировка, подбор мебели, проверка размеров, настройка света, рендеринг и
  подготовка спецификации;
  \item при работе с клиентом часто нужно показать несколько вариантов одного
  помещения, поэтому скорость генерации становится критичным фактором;
  \item автоматический пайплайн позволяет фиксировать правила, повторять
  эксперименты и сравнивать варианты по стоимости, наполнению и визуальному
  качеству;
  \item связка 3D-сцены с каталогом товаров повышает прикладную ценность
  результата, потому что пользователь получает не абстрактную картинку, а набор
  предметов, которые можно приобрести или заменить.
\end{enumerate}

В тексте главы важно показать, что задача не сводится только к генерации
красивого изображения. Необходимо получить структурированный результат:
расстановку объектов, список товаров, цены, изображения, файл сцены и отчет.
Именно это отличает прикладной пайплайн от демонстрационной генерации сцены.

\subsection{Сравнение существующих методов генерации планировок}

Для решения задачи были рассмотрены подходы, основанные на генерации
расположения объектов в помещении. В обзор следует включить M3D Layout,
DiffuScene, LEGO-Net и Infinigen. Для каждого подхода нужно кратко описать:
какие входные данные он принимает, какой результат формирует, насколько удобно
его запускать локально, можно ли управлять составом сцены и насколько легко
интегрировать результат с Blender и внешним каталогом товаров.

При описании M3D Layout, DiffuScene и LEGO-Net нужно сделать акцент не на том,
что эти методы плохие, а на том, что они оказались менее удобны для
практического пайплайна. Возможные ограничения:
\begin{enumerate}
  \item сложность воспроизведения окружения и запуска обученных моделей;
  \item зависимость от конкретных датасетов и форматов сцен;
  \item недостаточная управляемость результата при необходимости получить
  коммерчески применимую сцену;
  \item сложность последующей замены объектов на товары из каталога;
  \item отсутствие полного пайплайна от генерации до отчета и спецификации.
\end{enumerate}

\subsection{Обоснование выбора Infinigen}

Infinigen был выбран как базовый инструмент, поскольку он позволяет процедурно
создавать сцены, управлять параметрами генерации и получать результат в формате,
совместимом с Blender. Для данной работы это важно по нескольким причинам.
Во-первых, Blender используется как универсальная среда для загрузки,
модификации, рендеринга и экспорта 3D-сцен. Во-вторых, процедурная генерация
позволяет получать разные варианты сцены без ручного моделирования каждого
объекта. В-третьих, наличие полноценной 3D-сцены дает возможность выполнять
постобработку: добавлять отсутствующие элементы, менять материалы, заменять
объекты на магазинные аналоги и готовить отчет.

При этом Infinigen не решает всю задачу напрямую. В исходном виде результат
требует доработки: необходимо адаптировать расстановку под интерьерный сценарий,
контролировать состав объектов, улучшать освещение, извлекать описания
предметов и связывать элементы сцены с товарами из каталога. Поэтому в работе
Infinigen рассматривается не как готовое решение, а как базовый генератор,
вокруг которого построен дополнительный прикладной пайплайн.

\subsection{Постановка задачи}

На основе проведенного анализа задача работы формулируется следующим образом:
разработать пайплайн, который принимает описание помещения и набор требований к
интерьеру, генерирует 3D-сцену, улучшает расстановку объектов, подбирает
материалы и магазинные аналоги предметов, после чего формирует отчет с
визуализациями, ценами и выгружаемыми артефактами.

В конце первой главы нужно добавить краткий вывод: выбранный подход основан на
Infinigen, потому что он дает управляемую 3D-сцену в Blender, но для
практического применения требуется дополнительный слой обработки, описанный во
второй и третьей главах.

\section{Разработка пайплайна расстановки объектов}

\subsection{Общая архитектура решения}

Во второй главе нужно описать основной пайплайн расстановки. Рекомендуемая
логика изложения:
\begin{enumerate}
  \item входные данные: геометрия помещения, тип комнаты, ограничения,
  предпочтения пользователя, доступные каталоги;
  \item генерация базовой сцены с помощью Infinigen;
  \item анализ полученной сцены и извлечение списка объектов;
  \item улучшение расстановки и добавление недостающих предметов;
  \item настройка освещения и камер;
  \item замена синтетических объектов на товары из каталога;
  \item подготовка результата: рендеры, файл Blender, изображения товаров,
  цены и отчет.
\end{enumerate}

В этом месте стоит добавить схему пайплайна. В LaTeX ее можно вставить как
рисунок:
\begin{verbatim}
\begin{figure}[H]
  \centering
  \includegraphics[width=0.95\textwidth]{images/pipeline.png}
  \caption{Общая схема пайплайна генерации интерьерной сцены}
  \label{fig:pipeline}
\end{figure}
\end{verbatim}

\subsection{Улучшение результата расстановки Infinigen}

Здесь нужно описать, какие проблемы возникали в базовой расстановке и какие
правила были добавлены поверх Infinigen. Примеры проблем: некорректные
расстояния между объектами, отсутствие важных предметов для заданного типа
комнаты, неудачная ориентация мебели, пересечения объектов, недостаточно
выразительное освещение, несоответствие состава сцены пользовательскому запросу.

После этого нужно описать доработки:
\begin{enumerate}
  \item проверка состава сцены относительно ожидаемого набора объектов;
  \item добавление отсутствующих предметов;
  \item коррекция положения и ориентации объектов;
  \item настройка источников света;
  \item выбор камер и подготовка ракурсов для рендера;
  \item сохранение промежуточных и итоговых артефактов.
\end{enumerate}

\subsection{Получение описаний объектов и материалов}

Для сопоставления объектов сцены с товарами из каталога используются текстовые
описания. Их можно получать с помощью VLM-моделей, которые анализируют
изображение объекта или материала, и LLM-моделей, которые нормализуют описание,
выделяют ключевые признаки и приводят его к формату, удобному для поиска.

В тексте нужно указать, какие признаки извлекаются: тип объекта, стиль, цвет,
материал, форма, назначение, примерные размеры, совместимость с типом комнаты и
дополнительные ограничения. Для отделочных материалов полезно выделять тип
поверхности, цветовую гамму, фактуру, рисунок, область применения и визуальную
совместимость с остальной сценой.

\subsection{Замена синтетических объектов на магазинные аналоги}

После получения описаний выполняется поиск подходящих товаров в каталоге.
Синтетический объект из сцены сопоставляется с магазинным товаром по типу,
визуальным признакам, материалу, цвету, стилю и доступности 3D-модели. Если
точного совпадения нет, система подбирает несколько близких вариантов.

В этом разделе важно описать, что замена объекта должна учитывать не только
текстовое сходство, но и техническую пригодность модели: наличие 3D-файла,
совместимость формата, корректный масштаб, материалы и возможность загрузки в
Blender.

\subsection{Формирование отчета и выгрузка результата}

Итоговый результат пайплайна включает набор рендеров, файл сцены в формате
Blender, список использованных товаров, изображения товаров, цены и варианты
замен. Отчет нужен для того, чтобы результат можно было передать пользователю
или использовать на следующем этапе внедрения.

В конце второй главы нужно сделать вывод: разработанный пайплайн превращает
базовую генерацию сцены в прикладной результат, включающий визуализации,
магазинные объекты и структурированную выгрузку.

\section{Каталог материалов, замены и экспериментальная оценка}

\subsection{Сбор каталога предметов и материалов}

В третьей главе нужно описать, как формировался каталог товаров. Для каждого
товара желательно хранить название, категорию, изображения, ссылку на источник,
цену, параметры, описание, признаки стиля и путь к 3D-модели при наличии.

Отдельно следует раскрыть трудности, связанные с 3D-моделями. В практических
каталогах не всегда доступны файлы в нужном формате. Даже если модель есть, она
может быть представлена в формате, требующем конвертации, иметь некорректный
масштаб, неполные материалы или структуру, неудобную для автоматической загрузки.
Формат FBX выбран как один из наиболее практичных вариантов обмена между
разными визуализаторами, включая Blender и Unity, однако даже он не гарантирует
полную совместимость без проверки и нормализации.

\subsection{Описание предметов и отделочных материалов}

Как и во второй главе, здесь нужно описать использование VLM и LLM, но уже со
стороны каталога. Для товаров и материалов формируются нормализованные описания,
которые затем используются при поиске замен. Это позволяет сравнивать
синтетические объекты и реальные товары не только по названию категории, но и по
визуальным характеристикам.

Например, для кресла могут учитываться форма, цвет, материал обивки, стиль,
наличие деревянных или металлических элементов. Для напольного покрытия могут
учитываться оттенок, рисунок, тип поверхности, степень контраста и совместимость
с мебелью.

\subsection{Эксперименты и оценка качества}

Эксперименты должны показать, что предложенный пайплайн дает применимые
результаты. Рекомендуется описать несколько сценариев: разные типы комнат,
разные стилистические запросы, разные ограничения по составу объектов или
стоимости.

Оценку можно проводить по следующим критериям:
\begin{enumerate}
  \item корректность состава сцены;
  \item отсутствие грубых пересечений и визуальных ошибок;
  \item соответствие выбранных товаров исходным объектам;
  \item визуальная согласованность мебели и материалов;
  \item полнота итогового отчета;
  \item стоимость итогового набора товаров;
  \item наличие нескольких вариантов замен.
\end{enumerate}

Если используются оценки VLM, нужно явно описать, что модель выступает как
автоматический эксперт, который сравнивает изображение сцены или товара с
заданными критериями. Важно указать ограничения такой оценки: она не заменяет
экспертную проверку дизайнера, но позволяет быстро сравнивать большое количество
вариантов и находить явные ошибки.

\subsection{Расчет стоимости и варианты замен}

Для каждого ключевого объекта формируются варианты замены. Практически удобно
использовать три уровня:
\begin{enumerate}
  \item базовый вариант с минимальной стоимостью;
  \item основной вариант с наилучшим балансом цены и визуального соответствия;
  \item премиальный вариант с более высокой стоимостью или лучшим качеством.
\end{enumerate}

Такой подход позволяет использовать один и тот же пайплайн для разных бюджетов.
В отчете нужно показывать итоговую стоимость сцены, стоимость отдельных товаров
и ссылки на позиции каталога.

\subsection{Выводы по экспериментам}

В конце третьей главы нужно сформулировать, какие результаты были получены:
сколько сцен было сгенерировано, какие типы помещений проверялись, как часто
удавалось подобрать подходящие товары, какие ошибки оставались и какие
ограничения требуют дальнейшей доработки.

\section*{Заключение}
\addcontentsline{toc}{section}{Заключение}

В ходе работы был разработан пайплайн автоматизированной генерации интерьерных
3D-сцен, подбора предметов и отделочных материалов, а также формирования
итогового отчета с визуализациями и стоимостью. В качестве базового генератора
сцен был выбран Infinigen, поскольку он позволяет получать редактируемые
3D-сцены, совместимые с Blender, и служит удобной основой для дальнейшей
постобработки.

Были проанализированы существующие методы генерации планировок и 3D-сцен,
выявлены ограничения их применения в прикладном сценарии интерьерного дизайна.
Поверх базовой генерации был реализован дополнительный слой обработки:
улучшение расстановки, добавление недостающих объектов, настройка освещения,
описание объектов и материалов с использованием VLM и LLM, сопоставление с
каталогом товаров и подготовка выгружаемых артефактов.

Практический результат работы заключается в возможности получать не только
изображение интерьера, но и структурированный набор данных: сцену Blender,
рендеры, список товаров, изображения, цены и варианты замен. Такой формат
результата пригоден для дальнейшего внедрения и интеграции в рабочий процесс
компании.

В дальнейшем систему можно развивать за счет расширения каталога товаров,
улучшения поиска 3D-моделей, добавления экспертных правил для разных типов
помещений и включения ручной обратной связи от дизайнера или пользователя.

\newpage
\begin{thebibliography}{99}
\addcontentsline{toc}{section}{Список использованных источников}

\bibitem{example}
Автор И. О. Название источника. Город: Издательство, 2026.

\end{thebibliography}

\end{document}
